\GPTBegin{}

\section{Oriented geometry at a four-valent node}
\label{sec:calc1}

\subsection{The triple product and its geometric meaning}

Consider three classical vectors. Their scalar triple product
\begin{equation}
 d=\mathbf a\cdot(\mathbf b\times\mathbf c)
\end{equation}
measures an oriented volume: $|d|$ is the volume of the parallelepiped,
while its sign distinguishes its two handed orientations. Reversing all
three vectors changes that sign.

At a spin-network node, the angular momenta describe oriented face-area
vectors of a quantum polyhedron. For a tetrahedron these are four vectors
$\mathbf E_e$ satisfying $\sum_{e=1}^4\mathbf E_e=0$. Their triple product
has dimensions of volume squared, and the classical tetrahedral volume is
\begin{equation}
 V_{\rm tet}^2=\frac{2}{9}
 \left|\mathbf E_1\cdot(\mathbf E_2\times\mathbf E_3)\right|.
\end{equation}
See \cite{Bianchi2013QuantumGeometry} for the face-normal description.
The face-area vectors are distinct from the edge-vectors used to construct
a parallelepiped. This explains why a square root enters the volume formula.

We first calculate the dimensionless quantum triple product
\begin{equation}\label{eqn:q-def}
 \hat q=\mathbf J_1\cdot(\mathbf J_2\times\mathbf J_3)
 =\sum_{a,b,c=1}^3\varepsilon_{abc}J_1^aJ_2^bJ_3^c,
 \qquad J_e^a=\frac12\sigma_e^a.
\end{equation}
Here $e$ labels a leg, while $a,b,c$ label Cartesian components. The three
factors in each term act on different legs and commute, so $\hat q$ is
Hermitian. The relation to a chosen LQG volume prescription is stated below.

\subsection{The two physical intertwiner states}
\label{sub:cg}\label{sub:singlet-basis}

Take all four legs to be outgoing and to carry spin $j=\tfrac12$.
Gauge invariance requires
\begin{equation}
 (\mathbf J_1+\mathbf J_2+\mathbf J_3+\mathbf J_4)|\psi\rangle=0.
\end{equation}
Pairing legs $(12)$ and $(34)$ gives
\begin{equation}\label{eqn:cg-decomp}
 (\tfrac12\otimes\tfrac12)\otimes(\tfrac12\otimes\tfrac12)
 =(0\oplus1)\otimes(0\oplus1).
\end{equation}
A total singlet occurs once from $0\otimes0$ and once from $1\otimes1$.
The physical intertwiner space is therefore two-dimensional.


\subsubsection{Constructing the states from spinor polynomials}

A spinor has two components, written $z_e=(\alpha_e,\beta_e)$ for leg $e$.
For spin one-half we represent the basis states by the formal variables
$\alpha_e\leftrightarrow|0\rangle_e$ and
$\beta_e\leftrightarrow|1\rangle_e$. A polynomial with degree one in each
leg therefore represents a four-spin state: its monomials specify the
basis kets and their coefficients specify their amplitudes.

The elementary invariant is
\begin{equation}
 F_{ij}=\alpha_i\beta_j-\beta_i\alpha_j.
\end{equation}
Under a common $SU(2)$ rotation this antisymmetric product is unchanged,
because the rotation has determinant one. It represents the unnormalized
singlet $|01\rangle_{ij}-|10\rangle_{ij}$. This is the spinor-polynomial
construction used in the intertwiner description of
\cite{DittrichHnybida2016Ising}.

For four legs of spin one-half every spinor occurs once, giving three
pairings:
\begin{equation}
 A(z)=F_{12}F_{34},\qquad B(z)=F_{13}F_{24},\qquad C(z)=F_{14}F_{23}.
\end{equation}
Their expansion directly supplies the states. For instance,
\begin{align}
 A(z)&=(\alpha_1\beta_2-\beta_1\alpha_2)
       (\alpha_3\beta_4-\beta_3\alpha_4)\nonumber\\
 &\longleftrightarrow |0101\rangle-|0110\rangle-|1001\rangle+|1010\rangle.
\end{align}
Let $|A\rangle,|B\rangle,|C\rangle$ denote the normalized states represented
by $A(z)/2,B(z)/2,C(z)/2$. Expansion also gives
\begin{equation}
 A(z)-B(z)+C(z)=0,\qquad
 \langle A|B\rangle=\frac12.
\end{equation}
Thus the three pairings are linearly dependent. Since $A$ and $B$ are
independent, they span the two-dimensional intertwiner space. Orthogonalizing
$|B\rangle$ against $|A\rangle$ gives
\begin{equation}
 |\Phi_1\rangle=|A\rangle,\qquad
 |\Phi_2\rangle=\frac{2|B\rangle-|A\rangle}{\sqrt3}.
\end{equation}
In particular,
\begin{equation}
 |B\rangle=\frac12
 (|0011\rangle-|0110\rangle-|1001\rangle+|1100\rangle),
\end{equation}
so ordinary polynomial multiplication gives all coefficients in the
basis displayed below.

\subsubsection{The angular-momentum coupling explanation}

The same construction can be understood without spinor variables.
A pair of spin-one-half representations contains a singlet and a triplet.
The first four-spin singlet combines two pair-singlets; the second combines
two pair-triplets to total spin zero.

For a pair of spins, write
\begin{equation}
 |s\rangle=\frac{|01\rangle-|10\rangle}{\sqrt2},\qquad
 |t_+\rangle=|00\rangle,\quad
 |t_0\rangle=\frac{|01\rangle+|10\rangle}{\sqrt2},\quad
 |t_-\rangle=|11\rangle.
\end{equation}
An orthonormal basis is
\begin{align}
 |\Phi_1\rangle&=|s\rangle_{12}|s\rangle_{34}\nonumber\\
 &=\frac12(|0101\rangle-|0110\rangle-|1001\rangle+|1010\rangle),
 \label{eqn:Phi1}\\
 |\Phi_2\rangle&=\frac1{\sqrt3}
 (|t_+\rangle_{12}|t_-\rangle_{34}
 -|t_0\rangle_{12}|t_0\rangle_{34}
 +|t_-\rangle_{12}|t_+\rangle_{34})\nonumber\\
 &=\frac1{\sqrt{12}}
 (2|0011\rangle-|0101\rangle-|0110\rangle-|1001\rangle
 -|1010\rangle+2|1100\rangle).
 \label{eqn:Phi2}
\end{align}
The pair-singlet and pair-triplet sectors are orthogonal, and the displayed
coefficients give unit norm. This basis is also used in the CZX calculation below.

For a direct check, a general state with total magnetic number zero is
\begin{equation}
 |\psi\rangle=\alpha|0011\rangle+\beta|0101\rangle
 +\gamma|0110\rangle+\delta|1001\rangle
 +\epsilon|1010\rangle+\zeta|1100\rangle.
\end{equation}
The condition $S^+_{\rm tot}|\psi\rangle=0$ gives
\begin{equation}
 \alpha+\beta+\delta=0,\quad \alpha+\gamma+\epsilon=0,\quad
 \beta+\gamma+\zeta=0,\quad \delta+\epsilon+\zeta=0.
\end{equation}
These four equations are independent. Subtracting the second from the
first and the fourth from the third gives
\begin{equation}
 \beta+\delta-\gamma-\epsilon=0,\qquad
 \beta+\gamma-\delta-\epsilon=0.
\end{equation}
Adding and subtracting these relations gives $\epsilon=\beta$ and
$\delta=\gamma$. The first and third original equations then give
$\alpha=\zeta=-\beta-\gamma$. Only $\beta$ and $\gamma$ remain free.
This agrees with both the polynomial construction and angular-momentum addition.

Within the total-$M_z=0$ space, annihilation by $S^+_{\rm tot}$ is sufficient
for a singlet: the identity
$\mathbf S_{\rm tot}^2=S^-_{\rm tot}S^+_{\rm tot}
 +(S^z_{\rm tot})^2+S^z_{\rm tot}$ gives
$\mathbf S_{\rm tot}^2|\psi\rangle=0$.



\subsection{Choice of intertwiner basis}
\label{sub:manuscript-states}

We use the same orthonormal basis $\{|\Phi_1\rangle,|\Phi_2\rangle\}$ in
the main text and supplementary calculations. The first state couples
pairs $(12)$ and $(34)$ separately to spin zero; the second couples each
pair to spin one and then combines them to total spin zero.
The choice depends on the pairing of legs and the phases assigned to
the states. Another orthonormal basis describes the same intertwiner
space. Operator matrices and state coefficients change together, leaving
eigenvalues and physical predictions unchanged. With our leg ordering
and phase conventions, the orientation eigenstates will be
$|\pm\rangle=(|\Phi_1\rangle\pm i|\Phi_2\rangle)/\sqrt2$.

\subsection{Volume and signed orientation}
\label{sub:volume}

The triple product $\hat q$ is the common algebraic object needed here.
To distinguish it from the signed interior of a particular volume
prescription, write, on the four-valent singlet space,
\begin{equation}\label{eqn:Q-singlet}
 \hat Q_v=c_v\hat q,\qquad
 \hat V_v=\gamma^{3/2}\ell_P^3\sqrt{|\hat Q_v|}.
\end{equation}
The real dimensionless coefficient $c_v$ includes the normalization.
For the AL operator it also includes the edge-tangent orientation factors.
We leave it explicit rather than assign a numerical volume before these
choices are fixed.

For example, write the signed Ashtekar--Lewandowski interior as
\begin{equation}
 \hat Q_v=C\sum_{1\leq I<J<K\leq4}\epsilon_v(I,J,K)\hat q_{IJK},
 \qquad \hat q_{IJK}=\mathbf J_I\cdot(\mathbf J_J\times\mathbf J_K).
\end{equation}
Here $C$ is the chosen real normalization and $\epsilon_v(I,J,K)$ is the
sign of the determinant of the three outgoing edge tangents, or zero if
those tangents are linearly dependent. On gauge-invariant states, closure
and the angular-momentum commutators give
\begin{equation}
 \hat q_{124}=-\hat q_{123},\qquad
 \hat q_{134}=\hat q_{123},\qquad
 \hat q_{234}=-\hat q_{123}.
\end{equation}
Consequently
\begin{equation}
 c_v=C\,[\epsilon_v(1,2,3)-\epsilon_v(1,2,4)
          +\epsilon_v(1,3,4)-\epsilon_v(2,3,4)].
\end{equation}
This is a fixed-embedding calculation: these tangent signs are held fixed
when time reversal acts on the spin labels. If $c_v=0$, this volume
operator vanishes at the vertex. If $c_v\ne0$, its two signed eigenvalues
have opposite sign and its ordinary volume is degenerate. The
Rovelli--Smolin prescription places absolute values differently and should
be treated separately; no equality of the two prescriptions is assumed here.



\subsubsection{Matrix elements of the triple product}

There is a short way to evaluate $\hat q$ without expanding all sixteen
spin states. Define
\begin{equation}
 A=(\mathbf J_1+\mathbf J_2)^2,\qquad
 B=(\mathbf J_2+\mathbf J_3)^2.
\end{equation}
Since $[J_e^a,J_f^b]=i\delta_{ef}\sum_c\varepsilon_{abc}J_e^c$,
\begin{align}
 [A,B]&=4[\mathbf J_1\cdot\mathbf J_2,\mathbf J_2\cdot\mathbf J_3]
 \nonumber\\
 &=4i\sum_{a,b,c}\varepsilon_{abc}J_1^aJ_2^cJ_3^b
 =-4i\hat q.
\end{align}
Thus $\hat q=(i/4)[A,B]$, the identity given in
Eq.~(8) of Brunnemann and Thiemann \cite{Brunnemann2006Simplification}.

In the basis $(\Phi_1,\Phi_2)$ the pair $(12)$ has spin zero or one, so
\begin{equation}
 A=\begin{pmatrix}0&0\\0&2\end{pmatrix}.
\end{equation}
For spin one-half, the swap operator $P_{23}$ satisfies
$P_{23}=2\mathbf J_2\cdot\mathbf J_3+\tfrac12 I$, so
$B=I+P_{23}$. Here $I$ denotes the identity. The exchange operator is defined by
\begin{equation}
 P_{23}|m_1,m_2,m_3,m_4\rangle=|m_1,m_3,m_2,m_4\rangle.
\end{equation}
To see the identity, note that a swap has eigenvalue $+1$ on the symmetric
triplet and $-1$ on the antisymmetric singlet. From
\begin{equation}
 \mathbf J_2\cdot\mathbf J_3
 =\frac12\left[(\mathbf J_2+\mathbf J_3)^2-\frac32 I\right],
\end{equation}
its eigenvalues are $1/4$ on the triplet and $-3/4$ on the singlet.
The expression $2\mathbf J_2\cdot\mathbf J_3+\tfrac12 I$ therefore has
exactly the required swap eigenvalues. This standard spin-exchange identity
is also derived in \cite{YuSpinExchange}. Because
$\mathbf J_2^2=\mathbf J_3^2=\tfrac34 I$, it follows that $B=I+P_{23}$.
Swapping the middle two bits gives
\begin{equation}
 P_{23}|\Phi_1\rangle
 =\frac12(|0011\rangle-|0110\rangle-|1001\rangle+|1100\rangle).
\end{equation}
Taking its overlap with $\Phi_2$ gives
\begin{equation}
 \langle\Phi_2|P_{23}|\Phi_1\rangle
 =\frac{2+1+1+2}{2\sqrt{12}}=\frac{\sqrt3}{2}.
\end{equation}
The two diagonal overlaps are $1/2$ and $-1/2$, respectively. Hence
\begin{equation}
 B=\begin{pmatrix}3/2&\sqrt3/2\\\sqrt3/2&1/2\end{pmatrix},
 \qquad [A,B]=\begin{pmatrix}0&-\sqrt3\\\sqrt3&0\end{pmatrix},
\end{equation}
and therefore
\begin{equation}\label{eqn:q-matrix-result}
 \boxed{\hat q\big|_{\rm singlet}
 =\frac{\sqrt3}{4}\begin{pmatrix}0&-i\\i&0\end{pmatrix}.}
\end{equation}
The coefficient follows from $J^a=\sigma^a/2$ and the explicit definition
of $\hat q$; it contains no volume regularization constant.

\subsubsection{The two orientation states}

\begin{result}\label{res:volume}
 The states
 \begin{equation}
 |+\rangle=\frac{|\Phi_1\rangle+i|\Phi_2\rangle}{\sqrt2},\qquad
 |-\rangle=\frac{|\Phi_1\rangle-i|\Phi_2\rangle}{\sqrt2}
 \end{equation}
 satisfy $\hat q|\pm\rangle=\pm q_0|\pm\rangle$, where
 \begin{equation}\label{eqn:q0}
 q_0=\frac{\sqrt3}{4}.
 \end{equation}
 For $\hat Q_v=c_v\hat q$, their signed eigenvalues are $\pm c_vq_0$.
 Both have ordinary volume
 \begin{equation}\label{eqn:volume-eigenvalue}
 v_0=\gamma^{3/2}\ell_P^3\sqrt{|c_v|\,q_0}.
 \end{equation}
\end{result}
The labels $+$ and $-$ refer to the sign of the ordered triple product
$\hat q$. They agree with the signs of $\hat Q_v$ when $c_v>0$ and are
exchanged when $c_v<0$.

\subsubsection{Time reversal exchanges the orientations}

For one spin choose $\Theta_{1/2}=i\sigma_y\mathcal K$, with $\mathcal K$
complex conjugation in the computational basis. Then
$\Theta_{1/2}|0\rangle=-|1\rangle$ and
$\Theta_{1/2}|1\rangle=|0\rangle$. On four spins,
\begin{equation}
 \Theta=(i\sigma_y)^{\otimes4}\mathcal K,\qquad \Theta^2=1.
\end{equation}
Each displayed singlet is real and unchanged by this operation:
$\Theta|\Phi_1\rangle=|\Phi_1\rangle$ and
$\Theta|\Phi_2\rangle=|\Phi_2\rangle$.
Anti-linearity then gives
\begin{equation}
 \Theta|+\rangle=|-\rangle,\qquad\Theta|-\rangle=|+\rangle.
\end{equation}
Equivalently, in the orientation basis $\Theta=X\mathcal K$.
Since $\Theta\mathbf J_e\Theta^{-1}=-\mathbf J_e$,
\begin{equation}
 \Theta\hat q\Theta^{-1}=-\hat q,\qquad
 \Theta\hat V_v\Theta^{-1}=\hat V_v.
\end{equation}
Time reversal changes the sign of the oriented geometry while preserving
its ordinary volume. The further interpretation of this sign as a
spacetime time orientation is the physical proposal of the main text.

\GPTEnd{}
